\chapter{Related Work}
\label{sec:related}

%Common techniques for executing applications in the presence of failures
%include replication, checkpointing, and static or dynamic scheduling and
%resource provisioning.  In non-real-time settings, the objective is to
%optimize some performance metric~\cite{10.5555/2811302}.  In real-time
%settings, the objective is to optimize some timeliness metric given
%job/task deadlines~\cite{10.1145/3589950}.

%Examples in real-time
%settings, besides our motivating use case in Chapter~\ref{sec:use-case},
%include Adaptive Bitrate Streaming~\cite{7976298}, in which video
%resolution can be degraded to meet quality-of-service constraints, or the
%use of different precision levels in neural networks to meet
%deadlines~\cite{JIANG2020101775}.

This work targets scenarios in which an application task can implement
different trade-offs between execution time and output error, which is
often termed \emph{approximate computing}~\cite{7348659} and can be
achieved with hardware or software
techniques~\cite{10.1145/3716845,10.1145/3711683}.  Leveraging approximate
computing for real-time applications is a natural idea~\cite{259428,
Leung1995}, which has been explored for various workload and application
models~\cite{10.1002/cpe.4208, 10.1016/j.future.2019.02.019,
STAVRINIDES20101004, 585499, 10.1007/978-3-642-21713-5_20} and remains an
active area of research~\cite{10.1145/3589950}.  The scheduling of
real-time, approximate computing applications, for both aperiodic and
periodic scenarios, has been studied for transient errors, where the
probability of error depends on DVFS levels~\cite{10.1145/2442087.2442094,
10.1007/s11227-021-03764-x}, and for fail-stop failures~\cite{259428,
10.1177/0037549709340729, 10.1109/TC.2003.1183950, 7570524}.


A popular approximate computing model is the \emph{monotone}
model~\cite{259428}: each task consists of a mandatory portion that can
produce output with some error, followed by an optional portion that
incrementally reduces the error and can be interrupted at any time.  This
monotone model is not applicable to the use case that motivates this work,
as the level of approximation of a task must be decided at the onset of its
(re)execution, at which point its execution time and output error level are
fixed, and no usable intermediate output is produced.  Instead of the
monotone model, some works consider that each task comes with a primary
mode (costly, low-error) and an alternate mode (cheap, high-error).  For
instance, the authors in~\cite{10.1109/TC.2003.1183950} propose an approach
in which the primary is always scheduled before the alternate.  In this
work, we consider tasks that can have arbitrary numbers of execution modes.
A natural context in which tasks provide multiple execution modes is when
different DVFS levels lead to different trade-offs between a task's
execution time and the probability that its output is erroneous.
In~\cite{10.1145/2442087.2442094}, the authors focus on periodic tasks with
individual deadlines and seek to minimize energy consumption.
In~\cite{10.1007/s11227-021-03764-x}, the authors have a similar objective
but, like in this work, consider an application with an end-to-end
deadline. In both works, the objective is to minimize energy
consumption while guaranteeing a particular reliability, i.e., bounding the
probability that at least one task output is erroneous.  In contrast,
motivated by the use case in Chapter~\ref{sec:use-case}, in this work some
level of application output error is tolerable and we aim to minimize the
expected error.


The work most closely related to ours is~\cite{10678736}. It considers
deadline-constrained task chains subject to fail-stop failures, where task
execution modes offer different trade-offs between execution time and
output error.  The authors formulate a static resource provisioning
problem: determine optimal choices for the number of application replicas
to execute concurrently, the execution modes for each replica, and the
number of compute resources allocated to each replica, for a given resource
budget.  In contrast, we consider adaptive scheduling decisions throughout
execution, including replication, rollback/recovery, mode changes after
observed failures, reuse of persistent intermediate outputs, and
coordination across PEs. These additional mechanisms require reasoning
about the dynamic execution state rather than making a one-time resource
provisioning decision. To the best of our knowledge, the adaptive
scheduling problem in this work has not been studied previously.

